\chapter{pybind11 による C++ 拡張\mbox{モジュール}}
\label{app:pybind11}
\index{pybind11@pybind11|textbf}
\index{cpp extension@C++ 拡張モジュール|textbf}

Python のプロファイル結果から計算量の多い関数を特定した後、その関数だけを C++ で実装する方法がある。pybind11 は、C++ の関数やクラスを Python モジュールとして公開するためのライブラリである。本付録では、最小構成の C++ 拡張モジュールを作り、\code{uv pip install -e .} で開発用環境へインストールする手順を示す。

\section{C++ 拡張を選ぶ判断基準}

プロファイルによって少数の関数が処理時間を支配すると確認でき、NumPy などの既存実装では置き換えられない場合は、その関数を C++ で実装する候補となる。入出力、可視化、ワークフロー制御を Python に残し、計算量の多い関数だけを pybind11 で公開すれば、C++ との境界を限定できる。

具体例として、時刻列 \code{times} に対して、差の絶対値が \code{window} 以下となる全ペアの数を C++ で数える。この二重ループの計算量は、時刻数を \(N\) とすると \(O(N^2)\) である。

\section{pybind11 パッケージの最小構成}
\index{scikit-build-core@scikit-build-core|textbf}
\index{cmake@CMake|textbf}

ディレクトリ構成は次のようになる。

\begin{lstlisting}[numbers=none, xleftmargin=2\zw, title={fastcount\_cpp/ の例}]
pyproject.toml
CMakeLists.txt
cpp/
`-- bindings.cpp
src/
`-- fastcount_cpp/
    `-- __init__.py
\end{lstlisting}

\code{pyproject.toml} はビルド方法と Python パッケージの情報、\code{CMakeLists.txt} は C++ 拡張のコンパイル方法を記述する。\code{cpp/bindings.cpp} に C++ の処理と Python への公開方法を書き、\code{src/fastcount\_cpp/\_\_init\_\_.py} から利用者向けの関数を再公開する。以下では、各ファイルの役割と内容を構成順に説明する。

\subsection{\code{pyproject.toml}}

C++ 拡張には、Python パッケージの情報に加えてコンパイル方法を定めるビルドバックエンドが必要である。ここでは、\code{scikit-build-core}、pybind11、CMake を組み合わせる。

\begin{lstlisting}[numbers=none, xleftmargin=2\zw, title={pyproject.toml}]
[build-system]
requires = [
  "pybind11>=2.13",
  "scikit-build-core>=0.10",
]
build-backend = "scikit_build_core.build"

[project]
name = "fastcount_cpp"
version = "0.1.0"
requires-python = ">=3.12"
dependencies = ["numpy>=2.0"]

[tool.scikit-build]
wheel.packages = ["src/fastcount_cpp"]
\end{lstlisting}

\code{build-system} に書いたパッケージは、拡張モジュールのビルド環境で用いられる。インストール後の実行に必要なパッケージは \code{project.dependencies} に分けて記述する。

\subsection{\code{CMakeLists.txt}}
\index{cmakelists@\code{CMakeLists.txt}|textbf}

\code{CMakeLists.txt} には、使用する C++ 規格、コンパイル対象、インストール先を記述する。最小構成は次の通りである。

\begin{lstlisting}[numbers=none, xleftmargin=2\zw, title={CMakeLists.txt}]
cmake_minimum_required(VERSION 3.15)
project(fastcount_cpp LANGUAGES CXX)

find_package(pybind11 CONFIG REQUIRED)

pybind11_add_module(_core cpp/bindings.cpp)
target_compile_features(_core PRIVATE cxx_std_17)
install(TARGETS _core DESTINATION fastcount_cpp)
\end{lstlisting}

\code{pybind11\_add\_module} が、Python から import できる共有ライブラリを作る中心になる。ここで作るモジュール名を \code{\_core} にしておくと、Python 側では \code{fastcount\_cpp.\_core} として参照できる。

\subsection{C++ の実装}
\index{pybind11 module@\code{PYBIND11\_MODULE}|textbf}

\code{cpp/bindings.cpp} では、時刻差を数える C++ 関数を定義し、同じファイル内で pybind11 による公開名と引数名を指定する。

\begin{lstlisting}[language=C++, title={cpp/bindings.cpp}]
#include <cmath>
#include <cstddef>
#include <stdexcept>
#include <vector>

#include <pybind11/pybind11.h>
#include <pybind11/stl.h>

std::size_t count_pairs_within_window(
    const std::vector<double>& times,
    double window
) {
    if (window < 0.0) {
        throw std::invalid_argument("window must be non-negative");
    }

    std::size_t count = 0;
    const std::size_t n = times.size();

    for (std::size_t i = 0; i < n; ++i) {
        for (std::size_t j = i + 1; j < n; ++j) {
            if (std::abs(times[j] - times[i]) <= window) {
                ++count;
            }
        }
    }
    return count;
}

namespace py = pybind11;

PYBIND11_MODULE(_core, m) {
    m.doc() = "Count close pairs in C++.";
    m.def(
        "count_pairs_within_window",
        &count_pairs_within_window,
        py::arg("times"),
        py::arg("window"),
        "Return the number of pairs whose time difference is within window."
    );
}
\end{lstlisting}

\code{pybind11/stl.h} をインクルードすると、Python の list から \code{std::vector<double>} への変換が有効になる。したがって、この関数の呼び出し側に明示的な変換コードは不要である。負の時間窓は定義に反するため、C++ 側で \code{std::invalid\_argument} を送出する。

\subsection{Python から公開する API}

C++ 拡張を内部モジュール \code{\_core} として直接公開すると、利用者が実装上の名前を知る必要が生じる。そこで、パッケージの \code{\_\_init\_\_.py} から公開関数だけを再公開する。

\begin{lstlisting}[language=Python, title={src/fastcount\_cpp/\_\_init\_\_.py}]
"""Python interface for the C++ pair counter."""

from ._core import count_pairs_within_window

__all__ = ["count_pairs_within_window"]
\end{lstlisting}

この \code{\_\_init\_\_.py} に公開関数を再エクスポートすると、利用側は \code{from fastcount\_cpp import count\_pairs\_within\_window} と記述できる。C++ 側の内部モジュール名 \code{\_core} は、公開 API に現れない。

\section{拡張モジュールのビルドとインストール}
\index{editable install@editable install|textbf}

開発中は、プロジェクトの仮想環境へ editable install する。Python 側のファイルはインストール元を参照し、C++ ソースは再ビルドによってコンパイル済み拡張へ反映される。

\begin{lstlisting}[numbers=none, xleftmargin=2\zw]
$ uv sync
$ uv pip install -e .
$ uv run python -c "from fastcount_cpp import count_pairs_within_window; print(count_pairs_within_window([0.0, 0.1, 0.4], 0.25))"
\end{lstlisting}

最後の出力は次のようになる。

\begin{lstlisting}[numbers=none, xleftmargin=2\zw]
1
\end{lstlisting}

これは、\code{(0.0, 0.1)} だけが窓 \code{0.25} の中に入っていることを表している。

純粋な Python ファイルとは異なり、C++ ソースの変更はコンパイル済み拡張へ自動では反映されない。\code{bindings.cpp} を編集した後は、\code{uv pip install -e .} を再実行して拡張モジュールを再ビルドする。

\section{Python からの C++ 関数の呼び出し}

インストール後の拡張モジュールは、Python モジュールと同じ import 文で読み込める。

\begin{lstlisting}[language=Python]
from fastcount_cpp import count_pairs_within_window

times = [0.0, 0.1, 0.4, 0.55]
window = 0.2

count = count_pairs_within_window(times, window)
print(count)
\end{lstlisting}

Python へ公開する関数を限定し、その関数の引数、返り値、例外条件をテストで固定する。C++ への移植後もこの契約を保てば、利用側のコードを変更せずに内部実装を置き換えられる。

\section{C++ へ移す処理の選び方}

入出力、可視化、例外処理、引数処理まで C++ へ移すと、Python と C++ の境界が広がり、ビルドと保守が複雑になる。次のように、ワークフローを担う Python と計算コアを担う C++ の役割を分ける。

\begin{itemize}
\item Python: ファイルの読み込み、引数処理、図の保存、ワークフロー全体
\item C++: 処理時間を支配するループ、組み合わせ探索、反復数値計算
\end{itemize}

Python 側には引数検証を含む公開 API とワークフローを置き、C++ 側には測定上のボトルネックとなった計算コアを置く。この分離により、C++ のビルドが必要な範囲と、Python から検証する公開契約を限定できる。

ビルド設定と Python/C++ 間の型変換の詳細は、pybind11~\cite{pybind11_docs} と scikit-build-core~\cite{skbuild_docs} の公式資料を参照する。
